\documentclass[border=10pt, multi, tikz]{standalone}
\usepackage[T1]{fontenc}
\usepackage{import}
\subimport{./layers/}{init}
\usetikzlibrary{positioning}
\usetikzlibrary{3d}

% Paleta Okabe-Ito, importada de src/eda.py PALETTE
\definecolor{cdato}{HTML}{56B4E9}
\definecolor{bdato}{HTML}{A6DAF5}
\definecolor{cconv}{HTML}{0072B2}
\definecolor{bconv}{HTML}{56B4E9}
\definecolor{cpool}{HTML}{7F7F7F}
\definecolor{bpool}{HTML}{B0B0B0}
\definecolor{cgap}{HTML}{CC79A7}
\definecolor{bgap}{HTML}{E4B4CF}
\definecolor{cdensa}{HTML}{009E73}
\definecolor{bdensa}{HTML}{66C7AC}
\definecolor{csalida}{HTML}{D55E00}
\definecolor{bsalida}{HTML}{F0956B}
\definecolor{clatente}{HTML}{E69F00}
\definecolor{blatente}{HTML}{F3C86B}
\definecolor{cruido}{HTML}{F0E442}
\definecolor{bruido}{HTML}{F7F0A0}
\definecolor{crecurrente}{HTML}{004E7C}
\definecolor{brecurrente}{HTML}{56B4E9}

\begin{document}
\begin{tikzpicture}
\tikzstyle{connection}=[ultra thick,every node/.style={sloped,allow upside down},draw=\edgecolor,opacity=0.7]

\pic[shift={(0,0,0)}] at (0,0,0)
    {Box={
        name=mlplentrada,
        caption= ,
        xlabel={{" ","dummy"}},
        zlabel= ,
        fill=cdato,
        opacity=0.75,
        height=3.0,
        width=1.6,
        depth=13.91
        }
    };

\pic[shift={(2.15,0,0)}] at (mlplentrada-east)
    {RightBandedBox={
        name=mlpl0,
        caption= ,
        xlabel={{" ","dummy"}},
        zlabel= ,
        fill=cdensa,
        bandfill=bdensa,
        opacity=0.75,
        bandopacity=0.9,
        height=3.0,
        width=1.6,
        depth=23.03
        }
    };

\pic[shift={(2.15,0,0)}] at (mlpl0-east)
    {RightBandedBox={
        name=mlpl3,
        caption= ,
        xlabel={{" ","dummy"}},
        zlabel= ,
        fill=cdensa,
        bandfill=bdensa,
        opacity=0.75,
        bandopacity=0.9,
        height=3.0,
        width=1.6,
        depth=17.93
        }
    };

\pic[shift={(2.15,0,0)}] at (mlpl3-east)
    {RightBandedBox={
        name=mlpl6,
        caption= ,
        xlabel={{" ","dummy"}},
        zlabel= ,
        fill=cdensa,
        bandfill=bdensa,
        opacity=0.75,
        bandopacity=0.9,
        height=3.0,
        width=1.6,
        depth=14.2
        }
    };

\pic[shift={(2.15,0,0)}] at (mlpl6-east)
    {Box={
        name=mlpl9,
        caption= ,
        xlabel={{" ","dummy"}},
        zlabel= ,
        fill=csalida,
        opacity=0.75,
        height=3.0,
        width=1.6,
        depth=5.57
        }
    };
\draw [connection] (mlplentrada-east) -- node {\midarrow} (mlpl0-west);
\draw [connection] (mlpl0-east) -- node {\midarrow} (mlpl3-west);
\draw [connection] (mlpl3-east) -- node {\midarrow} (mlpl6-west);
\draw [connection] (mlpl6-east) -- node {\midarrow} (mlpl9-west);
\node[anchor=north, align=center, font=\scriptsize, text width=2.2cm, inner sep=2pt, fill=white, fill opacity=0.82, text opacity=1] at ([yshift=-26pt] mlplentrada-south) {\textbf{60 valores} \\ entrada};
\node[anchor=north, align=center, font=\scriptsize, text width=2.2cm, inner sep=2pt, fill=white, fill opacity=0.82, text opacity=1] at ([yshift=-26pt] mlpl0-south) {\textbf{256 unidades} \\ Linear 60->256 \\ ReLU \\ Dropout p=0.1};
\node[anchor=north, align=center, font=\scriptsize, text width=2.2cm, inner sep=2pt, fill=white, fill opacity=0.82, text opacity=1] at ([yshift=-26pt] mlpl3-south) {\textbf{128 unidades} \\ Linear 256->128 \\ ReLU \\ Dropout p=0.1};
\node[anchor=north, align=center, font=\scriptsize, text width=2.2cm, inner sep=2pt, fill=white, fill opacity=0.82, text opacity=1] at ([yshift=-26pt] mlpl6-south) {\textbf{64 unidades} \\ Linear 128->64 \\ ReLU \\ Dropout p=0.1};
\node[anchor=north, align=center, font=\scriptsize, text width=2.2cm, inner sep=2pt, fill=white, fill opacity=0.82, text opacity=1] at ([yshift=-26pt] mlpl9-south) {\textbf{1 unidad} \\ salida \\ Linear 64->1};
\node[anchor=south west, align=left, inner sep=0pt] at ([yshift=14pt] current bounding box.north west) {\begin{tabular}{@{}l@{}} \bfseries\large mlp\_l — MLP \\ \mdseries\small 56,833 parámetros · entrada (B, 60) · salida (B,) \end{tabular}};
\end{tikzpicture}
\end{document}
